% !TEX program = xelatex
\documentclass{resume}
\usepackage{zh_CN-Adobefonts_external} % 中文字体支持
\usepackage{linespacing_fix}
\usepackage{cite}
\usepackage{comment}

\usepackage[Chinese]{languageSelection}
\usepackage[color=blue]{notePlus}

% 所有间距由 resume.cls 统一控制

\begin{document}
\pagenumbering{gobble}

\CN{
  \name{黄锡超}
  \info{手机: (+86) 13736707307}{邮箱: hsn.zj@foxmail.com}{}{}
  \info{性别: 男}{微信: hsn7307}{在线简历: resume.huangsn.dev}{}
}

% =============================
% 教育背景
% =============================
\CN{
  \section{教育背景}
  \datedsubsection{\textbf{浙江理工大学科技与艺术学院}，计算机科学与技术（本科）}{2021.09 -- 2025.06}
}

% =============================
% 工作经历
% =============================
\CN{
  \section{工作经历}

  \datedsubsection{\textbf{武汉智识无垠科技有限公司}，后端开发工程师}{2025.06 -- 至今}
  \begin{itemize}[leftmargin=*]
    \item 参与 \textbf{核心平台} 后端架构演进，支撑 \textbf{十万级 DAU、千万级日解析量}。
    \item 负责\textbf{海外身份认证体系}建设，覆盖 Google OAuth / 邮箱双通道认证、CLI / 桌面端授权与订阅级多端会话限制。
    \item 负责\textbf{运营管理后台}全栈开发，覆盖公告、审批流、用户概览等模块，运营效率提升 \textbf{2 倍以上}。
    \item 负责 \textbf{B 端生态项目}开发（MCP 服务、pdfdeal Python 库等），降低用户接入门槛。
  \end{itemize}

  \datedsubsection{\textbf{杭州谐云科技有限公司}，后端开发工程师}{2024.07 -- 2025.04}
  \begin{itemize}[leftmargin=*]
    \item 承担 \textbf{平台功能开发与重构}，参与\textbf{产品评审与测试用例评审}，负责 \textbf{Jira L3 级工单}定位与处理。
    \item 推动容器网关由\textbf{命令式向声明式}架构转型，核心操作响应耗时 \textbf{20s → 300ms}。
  \end{itemize}
}

% =============================
% 项目经历
% =============================
\CN{
  \section{项目经历}

  \datedsubsection{\textbf{Doc2X — C/B 端智能文档解析平台}，后端开发}{2025.06 -- 至今}
  \begin{itemize}[leftmargin=*]
    \item \textbf{简介}：国内 TOP1 文档解析平台，PDF 转 Markdown / LaTeX / Word，日解析量 \textbf{千万级}，覆盖国内站与海外站双形态。
    \item \textbf{技术栈}：Go、gRPC、ConnectRPC、PostgreSQL、Redis、NATS、Kubernetes、OpenTelemetry
    \item \textbf{高并发治理}：高峰期解析、导出链路被免费长尾请求挤占，付费请求出现堆积告警。落地\textbf{分级限流 + 闲时排队}与多级缓存，将高峰流量平滑至闲时消化，\textbf{付费链路告警清零}。
    \item \textbf{通知系统}：短信、邮件、飞书 / 钉钉等通知散落于各业务模块，每新增渠道需改动业务代码。基于\textbf{策略 + 工厂模式}抽离多渠道通知中台，\textbf{新增渠道仅需扩展 Provider}，业务侧零改动。
    \item \textbf{身份认证}：海外用户默认以邮箱登录，对接 \textbf{Google OAuth} 打通注册与授权；为 CLI / 桌面端场景自建 \textbf{OAuth 2.0 授权码 + PKCE} 与刷新令牌轮换，整合后\textbf{重构登录体系}。
    \item \textbf{运营后台}：运营 / 客服原依赖工程师查数据并手工处理退款 / 发票 / 注销，沟通成本高且易漏配额回收。独立完成\textbf{公告、转账、用户概览、订单分析}等模块，并将\textbf{退款 / 发票 / 注销改造为自助流程}。
  \end{itemize}

  \datedsubsection{\textbf{观云台 — Kubernetes PaaS 平台}，后端开发}{2024.07 -- 2025.04}
  \begin{itemize}[leftmargin=*]
    \item \textbf{简介}：企业级 Kubernetes PaaS 平台，面向多租户提供容器服务与合规审计，支撑客户降低运维成本、提升交付效率。
    \item \textbf{技术栈}：Go、Gin、Kubernetes、Kubebuilder、MySQL
    \item \textbf{资源重构}：将资源管理服务由 \textbf{Java 重写为 Go}，并按 \textbf{K8s 原生 RBAC} 重整权限模型；期间通过\textbf{阅读 K8s 源码定位到一处上游校验缺陷}，并在资源层补齐预校验完成兜底。
    \item \textbf{架构升级}：容器网关原为命令式调用链，单次操作需顺序串接多个 API，\textbf{高峰期响应耗时 20s}且失败难以重试。引入 Reconcile 改造为\textbf{声明式 + 幂等收敛}，响应稳定收敛至 \textbf{300ms} 内。
    \item \textbf{合规巡检}：历史方案依赖运维人工核对配置基线，效率低且漏检率高。基于 CronJob + 自定义检测引擎实现\textbf{自动化基线审计}，纳管组件合规率达 \textbf{100\%}。
  \end{itemize}
}

% =============================
% 技能
% =============================
\CN{
  \section{专业技能}
  \begin{itemize}[leftmargin=*]
    \item \textbf{编程语言与框架}：熟练使用 Go；掌握 gRPC、ConnectRPC、Gin 及微服务架构设计。
    \item \textbf{云原生技术}：熟悉 Kubernetes 核心机制，具备通过源码分析定位问题的经验。
    \item \textbf{数据库与中间件}：熟练使用 PostgreSQL、Redis、NATS。
    \item \textbf{AI 与大模型}：熟练使用 Codex、Claude Code 等 AI 编程工具；具备 MCP 开发与开源经验。
  \end{itemize}
}

\end{document}
